\newpage

\section{মাজি}

\noindent অতীত কালের তারতম্যে মাজি ছয় রকম। কণা লাল, ফেল কালো।

\vspace{0.25cm}

\begingroup
\setlength{\arrayrulewidth}{0.5pt}
\setlength{\tabcolsep}{3pt}
\renewcommand{\arraystretch}{1.75}

\noindent\begin{tabularx}{\textwidth}{|C|C|C|}
  \hline
  \rowcolor{headergreen}
  {\color{white}\textbf{কাল}} &
  {\color{white}\textbf{রূপ}} &
  {\color{white}\textbf{অর্থ}} \\
  \hline
  \rowcolor{lightgreen}
  মুতলাক & {\arabicfont \fala} & সে করল \\
  \hline
  কারিব & {\arabicfont \madqarib} & সে এইমাত্র করল \\
  \hline
  \rowcolor{lightgreen}
  বাঈদ & {\arabicfont \madbaid} & সে করেছিল \\
  \hline
  ইস্তিমরারি & {\arabicfont \madist} & সে করতেছিল \\
  \hline
  \rowcolor{lightgreen}
  ইহতিমালি & {\arabicfont \madiht} & সম্ভবত সে করেছে \\
  \hline
  তামান্নাঈ & {\arabicfont \madtam} & যদি সে করত \\
  \hline
\end{tabularx}
\endgroup

\vspace{0.25cm}

\noindent মাজহুলে কণা একই থাকে, ফেল হয় {\arabicfont \majhulmadi}।
ইস্তিমরারিতে {\arabicfont \textarabic{\sfx{كَانَ}}} এর পর
{\arabicfont \majhulmud}।

\vspace{0.45cm}

% ─── Fil Madi Table ─────────────────────────────────────────
\begin{center}
  {\arabicfont \color{titlegreen} \textbf{\textarabic{الماضي المطلق المعروف}}}
\end{center}

\vspace{0.3cm}

\begin{center}
  {\arabicfont
    \textarabic{%
      \textbf{\color{titlegreen}علامة:}
      \enspace
      الجذر \textbf{\myroot{فَعَلَ}}
      \enspace|\enspace
      الضمير (علامة) \textbf{\color{suffixred}بالأحمر}%
    }
  }
\end{center}

\vspace{0.4cm}

\setlength{\arrayrulewidth}{0.5pt}
\setlength{\tabcolsep}{3pt}
\renewcommand{\arraystretch}{2.2}

\noindent\begin{tabularx}{\textwidth}{|C|C|C|}

  \hline
  \rowcolor{headergreen}
  {\arabicfont \color{white} \textbf{مُتَكَلِّم}} &
  {\arabicfont \color{white} \textbf{حَاضِر}}     &
  {\arabicfont \color{white} \textbf{غَائِب}}     \\
  \hline

  \rowcolor{lightgreen}
  {\arabicfont \faltu}   &
  {\arabicfont \falta}   &
  {\arabicfont \fala}    \\
  \hline

  {\arabicfont \falnaa}     &
  {\arabicfont \faltumaa}   &
  {\arabicfont \falaa}      \\
  \hline

  \rowcolor{lightgreen}
  &
  {\arabicfont \faltum}  &
  {\arabicfont \faluu}   \\
  \hline

  &
  {\arabicfont \falti}   &
  {\arabicfont \falat}   \\
  \hline

  \rowcolor{lightgreen}
  &
  {\arabicfont \faltumaaf} &
  {\arabicfont \falataa}   \\
  \hline

  &
  {\arabicfont \faltunna} &
  {\arabicfont \falna}    \\
  \hline

\end{tabularx}

\vspace{0.4cm}

\begin{center}
  {\arabicfont \color{gray} \small
    \textarabic{الجذر الأصلي بالأسود \textbullet\ العلامة (الضمير) بالأحمر}
  }
\end{center}

\vspace{0.3cm}

\begin{center}
  মাজহুল: {\arabicfont \majhulmadi}
  \quad ফাতে পেশ, আইনে যের, লামে যবর।
\end{center}
